\documentclass[11pt]{article}
\usepackage{amsmath,amssymb,amsthm,mathtools,enumitem}

\newtheorem{theorem}{Theorem}
\newtheorem{remark}[theorem]{Remark}

\begin{document}

\title{Scaling invariance of the angular lower-tail coefficient}
\author{}
\date{}
\maketitle

\section*{Statement of the problem}

Let \(C\) be a copula on \([0,1]^2\). For \(\gamma>0\) and \(\delta\in(0,2)\), define
\[
\kappa_L(\delta,\gamma;C)
=
\lim_{u\downarrow 0}
\frac{\log C\!\bigl(\gamma u^\delta,\gamma^{-1}u^{2-\delta}\bigr)}{\log u},
\]
whenever the limit exists.

We prove that, under suitable regularity conditions, \(\kappa_L(\delta,\gamma;C)\) does not depend on the choice of the scaling parameter \(\gamma\). We then verify that the same conclusion applies to the proposed copula \(C_{\rm PM}\).

\section*{Assumptions}

Fix \(\delta\in(0,2)\). Let \(\tilde\gamma>0\) be an arbitrary reference scaling parameter. We assume:

\begin{enumerate}[label=\textnormal{(A\arabic*)}]
\item \label{A1}
The quantity
\[
\tilde{\kappa}(d):=\kappa_L(d,\tilde\gamma;C)
=
\lim_{u\downarrow 0}
\frac{\log C\!\bigl(\tilde\gamma u^d,\tilde\gamma^{-1}u^{2-d}\bigr)}{\log u}
\]
exists for \(d\) in a neighborhood of the fixed point \(\delta\), and in particular
\(\tilde{\kappa}(\delta)\) is finite.

\item \label{A2}
The convergence above is locally uniform in \(d\) near \(\delta\): for every \(\varepsilon>0\), there exist \(r=r(\varepsilon)>0\) and \(u_0=u_0(\varepsilon)\in(0,1)\) such that
\[
\left|
\frac{\log C\!\bigl(\tilde\gamma u^d,\tilde\gamma^{-1}u^{2-d}\bigr)}{\log u}
-
\tilde{\kappa}(d)
\right|
\le \varepsilon
\]
for all \(u\in(0,u_0]\) and all \(d\in(0,2)\) satisfying \(|d-\delta|\le r\).

\item \label{A3}
The map \(d\mapsto \tilde{\kappa}(d)\) is continuous at \(d=\delta\).
\end{enumerate}

These are the regularity conditions needed for the \(\varepsilon\)-argument below.

\section*{Notation}

For \(\gamma>0\), \(d\in(0,2)\), and \(u\in(0,1)\), define
\[
E_\gamma(u;d)
:=
\frac{\log C\!\bigl(\gamma u^d,\gamma^{-1}u^{2-d}\bigr)}{\log u}.
\]
Then, by assumption,
\[
\lim_{u\downarrow 0} E_{\tilde\gamma}(u;d)=\tilde{\kappa}(d)
\]
for all \(d\) in a neighborhood of \(\delta\).

\section*{Main theorem}

\begin{theorem}\label{thm:gamma-invariance}
Assume \ref{A1}--\ref{A3}. Then, for every \(\gamma>0\),
\[
\kappa_L(\delta,\gamma;C)=\kappa_L(\delta,\tilde\gamma;C)=\tilde{\kappa}(\delta).
\]
Hence, for the fixed \(\delta\in(0,2)\), the value of \(\kappa_L(\delta,\gamma;C)\) is independent of \(\gamma\).
\end{theorem}

\begin{proof}
Fix \(\gamma>0\), and set
\[
\Delta:=\log\!\left(\frac{\gamma}{\tilde\gamma}\right).
\]
For each \(u\in(0,1)\), let
\[
t:=\log u<0,
\qquad
\delta_u:=\delta+\frac{\Delta}{t}.
\]
We first prove the key identity
\[
E_\gamma(u;\delta)=E_{\tilde\gamma}(u;\delta_u).
\]

Indeed, since \(u=e^t\) and \(e^\Delta=\gamma/\tilde\gamma\),
\[
\tilde\gamma u^{\delta_u}
=
\tilde\gamma e^{(\delta+\Delta/t)t}
=
\tilde\gamma e^{\delta t}e^\Delta
=
\gamma u^\delta,
\]
and similarly
\[
\tilde\gamma^{-1}u^{2-\delta_u}
=
\tilde\gamma^{-1}e^{(2-\delta-\Delta/t)t}
=
\tilde\gamma^{-1}e^{(2-\delta)t}e^{-\Delta}
=
\gamma^{-1}u^{2-\delta}.
\]
Substituting these identities into the definition of \(E_{\tilde\gamma}\) yields
\[
E_{\tilde\gamma}(u;\delta_u)
=
\frac{\log C\!\bigl(\tilde\gamma u^{\delta_u},\tilde\gamma^{-1}u^{2-\delta_u}\bigr)}{\log u}
=
\frac{\log C\!\bigl(\gamma u^\delta,\gamma^{-1}u^{2-\delta}\bigr)}{\log u}
=
E_\gamma(u;\delta).
\]

We now prove that
\[
\lim_{u\downarrow 0}E_\gamma(u;\delta)=\tilde{\kappa}(\delta).
\]
Let \(\varepsilon>0\) be arbitrary. We must show that there exists \(u^*\in(0,1)\) such that
\[
\bigl|E_\gamma(u;\delta)-\tilde{\kappa}(\delta)\bigr|\le \varepsilon
\qquad\text{for all }u\in(0,u^*].
\]
Using the identity just established, it is enough to control
\[
\bigl|E_{\tilde\gamma}(u;\delta_u)-\tilde{\kappa}(\delta)\bigr|.
\]
Add and subtract \(\tilde{\kappa}(\delta_u)\), and apply the triangle inequality:
\begin{equation}\label{eq:triangle}
\bigl|E_{\tilde\gamma}(u;\delta_u)-\tilde{\kappa}(\delta)\bigr|
\le
\bigl|E_{\tilde\gamma}(u;\delta_u)-\tilde{\kappa}(\delta_u)\bigr|
+
\bigl|\tilde{\kappa}(\delta_u)-\tilde{\kappa}(\delta)\bigr|.
\end{equation}
Alternatively, if one wishes to compare directly with \(E_{\tilde\gamma}(u;\delta)\), then one may write
\begin{align}
\bigl|E_{\tilde\gamma}(u;\delta_u)-E_{\tilde\gamma}(u;\delta)\bigr|
&\le
\bigl|E_{\tilde\gamma}(u;\delta_u)-\tilde{\kappa}(\delta_u)\bigr|
+
\bigl|\tilde{\kappa}(\delta_u)-\tilde{\kappa}(\delta)\bigr|
+
\bigl|\tilde{\kappa}(\delta)-E_{\tilde\gamma}(u;\delta)\bigr|.
\label{eq:full-triangle}
\end{align}
We now control the relevant terms.

\medskip
\noindent
\textbf{Step 1: control of \(\delta_u-\delta\).}
Since \(\Delta\) is fixed and \(t=\log u\to-\infty\) as \(u\downarrow 0\),
\[
\delta_u-\delta=\frac{\Delta}{\log u}\longrightarrow 0.
\]
Therefore, for every \(r>0\), there exists \(u_1\in(0,1)\) such that
\[
|\delta_u-\delta|\le r
\qquad\text{for all }u\in(0,u_1].
\]

\medskip
\noindent
\textbf{Step 2: control of \(\bigl|\tilde{\kappa}(\delta_u)-\tilde{\kappa}(\delta)\bigr|\).}
By continuity of \(\tilde{\kappa}\) at \(\delta\), there exists \(r=r(\varepsilon)>0\) such that
\[
|d-\delta|\le r
\quad\Longrightarrow\quad
|\tilde{\kappa}(d)-\tilde{\kappa}(\delta)|\le \frac{\varepsilon}{2}.
\]
By Step 1, there exists \(u_2\in(0,1)\) such that
\[
|\delta_u-\delta|\le r
\qquad\text{for all }u\in(0,u_2].
\]
Hence,
\begin{equation}\label{eq:kappa-cont}
|\tilde{\kappa}(\delta_u)-\tilde{\kappa}(\delta)|\le \frac{\varepsilon}{2}
\qquad\text{for all }u\in(0,u_2].
\end{equation}

\medskip
\noindent
\textbf{Step 3: control of \(\bigl|E_{\tilde\gamma}(u;\delta_u)-\tilde{\kappa}(\delta_u)\bigr|\).}
By local uniform convergence, there exist \(r'=r'(\varepsilon)>0\) and \(u_3=u_3(\varepsilon)\in(0,1)\) such that
\[
\left|
E_{\tilde\gamma}(u;d)-\tilde{\kappa}(d)
\right|
\le \frac{\varepsilon}{2}
\]
for all \(u\in(0,u_3]\) and all \(d\) satisfying \(|d-\delta|\le r'\). Again by Step 1, there exists \(u_4\in(0,1)\) such that
\[
|\delta_u-\delta|\le r'
\qquad\text{for all }u\in(0,u_4].
\]
Therefore,
\begin{equation}\label{eq:uniform}
\left|
E_{\tilde\gamma}(u;\delta_u)-\tilde{\kappa}(\delta_u)
\right|
\le \frac{\varepsilon}{2}
\qquad\text{for all }u\in(0,\min\{u_3,u_4\}].
\end{equation}

\medskip
\noindent
\textbf{Step 4: conclusion.}
Let
\[
u^*:=\min\{u_2,u_3,u_4\}.
\]
Then for all \(u\in(0,u^*]\), combining \eqref{eq:triangle}, \eqref{eq:kappa-cont}, and \eqref{eq:uniform}, we obtain
\[
\bigl|E_{\tilde\gamma}(u;\delta_u)-\tilde{\kappa}(\delta)\bigr|
\le
\frac{\varepsilon}{2}+\frac{\varepsilon}{2}
=
\varepsilon.
\]
Since \(E_\gamma(u;\delta)=E_{\tilde\gamma}(u;\delta_u)\), it follows that
\[
\bigl|E_\gamma(u;\delta)-\tilde{\kappa}(\delta)\bigr|
\le \varepsilon
\qquad\text{for all }u\in(0,u^*].
\]
This proves that
\[
\lim_{u\downarrow 0}E_\gamma(u;\delta)=\tilde{\kappa}(\delta).
\]
Equivalently,
\[
\kappa_L(\delta,\gamma;C)=\tilde{\kappa}(\delta)=\kappa_L(\delta,\tilde\gamma;C).
\]
The proof is complete.
\end{proof}

\begin{remark}
If one prefers to prove directly that
\[
\lim_{u\downarrow 0}\bigl(E_\gamma(u;\delta)-E_{\tilde\gamma}(u;\delta)\bigr)=0,
\]
then one may use \eqref{eq:full-triangle} instead of \eqref{eq:triangle}. The three terms on the right-hand side are handled respectively by local uniform convergence, continuity of \(\tilde{\kappa}\), and ordinary convergence of \(E_{\tilde\gamma}(u;\delta)\to \tilde{\kappa}(\delta)\).
\end{remark}

\section*{Application to the Pareto-mixture copula}

We now state the corresponding conclusion for the proposed copula \(C_{\rm PM}\).

\begin{theorem}\label{thm:pm}
Suppose that for the copula \(C_{\rm PM}\), the limit
\[
\kappa_L(\delta,\tilde\gamma;C_{\rm PM})
=
\lim_{u\downarrow 0}
\frac{\log C_{\rm PM}\!\bigl(\tilde\gamma u^\delta,\tilde\gamma^{-1}u^{2-\delta}\bigr)}{\log u}
\]
exists for \(d\) in a neighborhood of \(\delta\), that the convergence is locally uniform in \(d\), and that the resulting function \(d\mapsto \kappa_L(d,\tilde\gamma;C_{\rm PM})\) is continuous at \(d=\delta\). Then, for every \(\gamma>0\),
\[
\kappa_L(\delta,\gamma;C_{\rm PM})
=
\kappa_L(\delta,\tilde\gamma;C_{\rm PM}).
\]
Hence \(\kappa_L(\delta,\gamma;C_{\rm PM})\) is independent of \(\gamma\).
\end{theorem}

\begin{proof}
The proof is exactly the same as that of Theorem \ref{thm:gamma-invariance}, with \(C\) replaced everywhere by \(C_{\rm PM}\). Indeed, define
\[
E_\gamma^{\rm PM}(u;d)
:=
\frac{\log C_{\rm PM}\!\bigl(\gamma u^d,\gamma^{-1}u^{2-d}\bigr)}{\log u}.
\]
For fixed \(\gamma>0\), let
\[
\Delta=\log\!\left(\frac{\gamma}{\tilde\gamma}\right),
\qquad
\delta_u=\delta+\frac{\Delta}{\log u}.
\]
Then the same algebra gives
\[
E_\gamma^{\rm PM}(u;\delta)=E_{\tilde\gamma}^{\rm PM}(u;\delta_u).
\]
The same triangle inequality argument, together with local uniform convergence and continuity, yields
\[
\lim_{u\downarrow 0}E_\gamma^{\rm PM}(u;\delta)
=
\kappa_L(\delta,\tilde\gamma;C_{\rm PM}).
\]
Therefore,
\[
\kappa_L(\delta,\gamma;C_{\rm PM})
=
\kappa_L(\delta,\tilde\gamma;C_{\rm PM}),
\]
as claimed.
\end{proof}

\section*{Conclusion}

Under the stated regularity conditions, the angular lower-tail coefficient
\[
\kappa_L(\delta,\gamma;C)
=
\lim_{u\downarrow 0}
\frac{\log C\!\bigl(\gamma u^\delta,\gamma^{-1}u^{2-\delta}\bigr)}{\log u}
\]
does not depend on the scaling parameter \(\gamma\). The same conclusion holds for the proposed Pareto-mixture copula \(C_{\rm PM}\), provided the same existence, local uniformity, and continuity assumptions are satisfied.

\end{document}